%% kapa: exam=1 round=24 year=2013 period=1 form=A subject=civil_law subject_ko=민법 seq=7 number=7
%% source: Q-Net 감정평가사 기출문제 (cst00309) artlSeq=5073802 file=extracted/1008/5073802/2013년 제24회 감정평가사 1차시험 기출문제/1교시(A형).hwp
%% conversion: kordoc 4.24.51 -> tools/gen1.py (auto); review=proofread-sonnet; flags=none
\begin{question}{7}
주물․종물에 관한 설명으로 옳은 것은? (다툼이 있으면 판례에 의함)
\begin{choices}
  \item 어떤 물건이 주물의 소유자나 이용자의 상용에 공여되고 있다면 주물의 효용과 직접 관계없어도 종물이다.
  \item 주유소의 지하에 매설된 유류저장탱크는 주유소의 종물이다.
  \item 종물은 주물의 처분에 따라야 하므로 종물만의 처분을 당사자가 임의로 정할 수는 없다.
  \item 건물에 대한 저당권의 효력은 그 건물소유를 목적으로 하는 지상권에도 미친다.
  \item 법률의 규정 또는 다른 약정이 없으면 주물에 설정된 저당권은 저당권설정 이후의 종물에 그 효력이 미치지 않는다.
\end{choices}
\end{question}
\answer{4}
